\documentclass[a4paper]{arrowhead}

\usepackage[yyyymmdd]{datetime}
\usepackage{etoolbox}
\usepackage[T1]{fontenc}
\usepackage[utf8]{inputenc}
\usepackage{multirow}
\usepackage{longtable}
\usepackage{amssymb}
\usepackage{float}
\usepackage{tikz}
\usetikzlibrary{positioning,arrows.meta,fit,calc,shapes.misc,shapes.geometric}

\renewcommand{\dateseparator}{-}
\setlength{\emergencystretch}{3em}
\lstset{basicstyle=\ttfamily\footnotesize,breaklines=true,columns=fullflexible,frame=single,rulecolor=\color{gray!50},backgroundcolor=\color{gray!6}}
\tikzset{
  ahsys/.style={draw=ArrowheadBlue,very thick,fill=ArrowheadBlue!6,rounded corners=2pt,minimum width=3.2cm,minimum height=1.1cm,align=center,font=\small\bfseries},
  ahext/.style={draw=ArrowheadGrey,thick,dashed,fill=gray!5,rounded corners=2pt,minimum width=3.2cm,minimum height=1.0cm,align=center,font=\small},
  ahsrv/.style={draw=ArrowheadPurple,thick,fill=ArrowheadPurple!8,rounded corners=6pt,minimum height=0.55cm,inner xsep=4pt,align=center,font=\scriptsize\ttfamily},
  ahcsm/.style={draw=ArrowheadTeal,thick,fill=ArrowheadTeal!8,rounded corners=6pt,minimum height=0.55cm,inner xsep=4pt,align=center,font=\scriptsize\ttfamily},
  ahdb/.style={draw=ArrowheadOrange,thick,fill=ArrowheadOrange!8,cylinder,shape border rotate=90,aspect=0.25,minimum height=0.9cm,align=center,font=\scriptsize},
  ahflow/.style={-{Stealth[length=2.2mm]},thick,ArrowheadGrey},
  ahlabel/.style={font=\scriptsize,fill=white,inner sep=1pt,align=center},
}

\newcommand{\fparam}[1]{\textit{\textcolor{ArrowheadBlue}{#1}}}
%% Special references
\newcommand{\mref}[1]{{\textcolor{ArrowheadPurple}{\hyperref[sec:model:#1]{#1}}}}
\newcommand{\oref}[1]{{\textcolor{ArrowheadBlue}{\hyperref[sec:operations:#1]{#1}}}}
\newcommand{\pdef}[1]{{\textcolor{ArrowheadGrey}{#1\label{sec:model:primitives:#1}\label{sec:model:primitives:#1s}}}}
\newcommand{\pref}[1]{{\textcolor{ArrowheadGrey}{\hyperref[sec:model:primitives:#1]{#1}}}}

\newrobustcmd\osubsection[5]{
  \addtocounter{subsection}{1}
  \addcontentsline{toc}{subsection}{\protect\numberline{\thesubsection}operation \textcolor{ArrowheadBlue}{#1}}
  \subsection*{
    \thesubsection\quad
    #2 \textcolor{ArrowheadPurple}{#3} \\
    \small
    \hspace*{0.075\textwidth}\begin{minipage}{0.1\textwidth}
      \vspace*{1mm}
      Operation: \\
      \notblank{#4}{Input: \\}{}
      \notblank{#5}{Output: \\}{}
    \end{minipage}
    \begin{minipage}{0.825\textwidth}
      \vspace*{1mm}
      \textcolor{ArrowheadBlue}{#1} \\
      \notblank{#4}{\mref{#4} \\}{}
      \notblank{#5}{\mref{#5} \\}{}
    \end{minipage}
  }
  \label{sec:operations:#1}
}
\newrobustcmd\msubsection[2]{
  \addtocounter{subsection}{1}
  \addcontentsline{toc}{subsection}{\protect\numberline{\thesubsection}#1 \textcolor{ArrowheadPurple}{#2}}
  \subsection*{\thesubsection\quad#1 \textcolor{ArrowheadPurple}{#2}}
  \label{sec:model:#2} \label{sec:model:#2s}
}

\begin{document}

%% Arrowhead Document Properties
\ArrowheadTitle{sysMonitor AMQP/TLS/JSON}
\ArrowheadServiceID{sysMonitor}
\ArrowheadType{Interface Design Description}
\ArrowheadTypeShort{IDD}
\ArrowheadVersion{5.4.0}
\ArrowheadDate{2026-11-25}
\ArrowheadAuthor{Documentation maintenance}
\ArrowheadStatus{REVIEW DRAFT}
\ArrowheadContact{Repository maintainers}
\ArrowheadFooter{\href{www.arrowhead.eu}{www.arrowhead.eu}}
\ArrowheadSetup
%%

%% Front Page
\begin{center}
  \vspace*{1cm}
  \huge{\arrowtitle}

  \vspace*{0.2cm}
  \LARGE{\arrowtype}
  \vspace*{1cm}

  \Large{Service ID: \textit{``\arrowid''}}
  \vspace*{\fill}

  

  \vspace*{1cm}
  \vspace*{\fill}

  % Front Page Abstract
  \begin{abstract}
    This document describes the AMQP/TLS/JSON interface of the WorkflowManager's own observation types of the \texttt{sysMonitor} service, service version 1.0.0, documentation edition 5.4.0. It realises the \texttt{sysMonitor} SD of the WorkflowManager.
  \end{abstract}

  \vspace*{1cm}
\end{center}

\newpage
%%

%% Table of Contents
\tableofcontents
\newpage
%%

\section{Overview}
\label{sec:overview}
The WorkflowManager publishes its observations as persistent AMQP 0-9-1 messages with content type
\texttt{application/\allowbreak{}json}~\cite{rfc8259} to the topic exchange of the local cloud (\texttt{sysmonitor.\allowbreak{}exchange}) and waits
for the broker's confirmation. The contract of the envelope and of \texttt{echo}, \texttt{power} and \texttt{status} is
SysMonitor-AMQP-IDD-5.2.2; this document adds the types \texttt{executor\allowbreak{}Liveness} and \texttt{workItemLost}.

Section~\ref{sec:operations} specifies the publications, Section~\ref{sec:model} the data models,
Section~\ref{sec:subscription} the subscription, Section~\ref{sec:security} the security requirements and
Section~\ref{sec:implementations} the implementations.

\newpage
\section{Service Operations}
\label{sec:operations}

{\small
\begin{longtable}{|p{3.2cm}|p{1.5cm}|p{7.6cm}|p{3.5cm}|}

\hline
\rowcolor{gray!33} Type & Method & Routing key & Value \\ \hline
\endhead
\oref{executorLiveness} & publish & \texttt{\{cloud\}.sysmonitor.\allowbreak{}WorkflowManager.\allowbreak{}\{instance\}.\allowbreak{}executorLiveness} & \mref{ExecutorLiveness} \\ \hline
\oref{workItemLost} & publish & \texttt{\{cloud\}.sysmonitor.\allowbreak{}WorkflowManager.\allowbreak{}\{instance\}.\allowbreak{}workItemLost} & \mref{WorkItemLost} \\ \hline
\end{longtable}
}

In the headings \texttt{\{base\}} is \texttt{\{cloud\}.sysmonitor.WorkflowManager.\{instance\}}.
\osubsection{executorLiveness}{publish}{\{base\}.executorLiveness}{}{ExecutorLiveness}
State. It is published in every publication cycle (\texttt{sysmonitor.\allowbreak{}publish.\allowbreak{}interval}) and at once when an
executor becomes suspected, is cleared, refuses the question, or a work item is lost. A subscriber therefore learns of a
possible loss at the first failed check, before the grace period \texttt{lost\allowbreak{}After\allowbreak{}Seconds} has passed, and can secure the
equipment. Quality is \texttt{good} in the states \texttt{ok} and \texttt{disabled}, otherwise \texttt{uncertain};
\texttt{heldWorkItems} is as of the last check.
\osubsection{workItemLost}{publish}{\{base\}.workItemLost}{}{WorkItemLost}
Event, one per lost work item, with quality \texttt{uncertain}. It is kept, and published again with the same
observation id, until the broker has confirmed it, so a subscriber may receive it more than once. At most 1000
unpublished events are kept, in memory: a restart of the WorkflowManager before the broker confirms an event loses the
event (the worklog of \texttt{workflow\allowbreak{}Management} still holds the loss).

\subsection{Envelope}
Every message carries the common envelope:

{\small
\begin{longtable}{|p{5.6cm}|p{10.2cm}|}

\hline
\rowcolor{gray!33} Field & Value \\ \hline
\endhead
\texttt{schemaVersion} & \texttt{1.0} \\ \hline
\texttt{observationId} & \texttt{WorkflowManager-\{type\}-\{UUID\}}; repeated unchanged when an event is published again. \\ \hline
\texttt{provider\allowbreak{}System}, \texttt{service\allowbreak{}Instance}, \texttt{observation\allowbreak{}Type} & As in the routing key; SysMonitor refuses a difference. \\ \hline
\texttt{sourceTime} & When the value was taken (UTC). \\ \hline
\texttt{value} & \mref{ExecutorLiveness} or \mref{WorkItemLost}. \\ \hline
\texttt{unit} & \texttt{1} for \texttt{executor\allowbreak{}Liveness}, null for \texttt{workItemLost}. \\ \hline
\texttt{quality} & \texttt{good} or \texttt{uncertain}. \\ \hline
\texttt{sequence}, \texttt{metadata} & As for the common observation types. \\ \hline
\end{longtable}
}


\begin{lstlisting}[caption={A lost work item (shortened, from the test with a killed executor)}]
{ "schemaVersion": "1.0",
  "observationId": "WorkflowManager-workItemLost-ef1abb7c-4d77-433d-8a7d-0c4aa8a78041",
  "providerSystem": "WorkflowManager", "serviceInstance": "WorkflowManager",
  "observationType": "workItemLost", "sourceTime": "2026-11-25T08:23:29.607Z",
  "value": { "workItemId": "8867b471-47c0-4392-abd9-38f949345e2b", "sessionId": 1,
    "action": "a", "step": "work", "attempt": 1, "bound": true,
    "recipe": "rebind@1.0.0", "executor": "WfeA", "cause": "unreachable",
    "detail": "executor WfeA cannot be reached: ...",
    "lastSeen": "2026-11-25T08:22:42.495Z",
    "suspectedSince": "2026-11-25T08:23:17.472Z",
    "lostAt": "2026-11-25T08:23:29.607Z", "detectionSeconds": 47.112,
    "equipmentMayBeInUse": true, "outcome": "rebind", "workRepeated": true },
  "unit": null, "quality": "uncertain" }
\end{lstlisting}


\section{Subscription}
\label{sec:subscription}
SysMonitor offers the types it finds in the registration. A subscription to both, at SysMonitor's management service:

\begin{lstlisting}
PUT /sysmonitor/management/subscriptions/WorkflowManager
{ "observationTypes": ["executorLiveness", "workItemLost"], "enabled": true }
\end{lstlisting}

SysMonitor then publishes them on MQTT as \texttt{arrowhead/\{cloud\}/sysmonitor/\allowbreak{}data/workItemLost/\allowbreak{}WorkflowManager} and
\texttt{.\allowbreak{}.\allowbreak{}.\allowbreak{}/\allowbreak{}executor\allowbreak{}Liveness/\allowbreak{}Workflow\allowbreak{}Manager}. A message that arrives twice is stored once (\texttt{observationId}).

\newpage
\section{Data Models}
\label{sec:model}
\msubsection{struct}{ExecutorLiveness}
Which executors are suspected of no longer running the work items they hold.

{\small
\begin{longtable}{|p{3.7cm}|p{4.2cm}|p{1.6cm}|p{6.3cm}|}

\hline
\rowcolor{gray!33} Field & Type & Mandatory & Description \\ \hline
\endhead
\texttt{state} & \mref{LivenessState} & yes & Summary of the last check. \\ \hline
\texttt{heldWorkItems} & \pref{Integer} & yes & Work items dispatched to, claimed by or bound to executors. \\ \hline
\texttt{suspected\allowbreak{}Executors} & \pref{List}$<$\mref{SuspectedExecutor}$>$ & yes & Executors that did not confirm a work item they hold. \\ \hline
\texttt{unchecked\allowbreak{}Executors} & \pref{List}$<$\pref{SystemName}$>$ & yes & Executors that refuse \texttt{workflow\allowbreak{}Execution.\allowbreak{}get-\allowbreak{}status}; nothing is known about them. \\ \hline
\texttt{lostWorkItems} & \pref{Integer} & yes & Work items lost since the WorkflowManager started. \\ \hline
\texttt{checkSeconds} & \pref{Duration} & yes & \texttt{workflow.\allowbreak{}executor.\allowbreak{}liveness.\allowbreak{}check-\allowbreak{}seconds}. \\ \hline
\texttt{lost\allowbreak{}After\allowbreak{}Seconds} & \pref{Duration} & yes & \texttt{workflow.\allowbreak{}executor.\allowbreak{}liveness.\allowbreak{}lost-\allowbreak{}after-\allowbreak{}seconds}. \\ \hline
\end{longtable}
}

\msubsection{enum}{LivenessState}
Summary of the liveness check.
 Encoded as a JSON string.
{\small
\begin{longtable}{|p{3.9cm}|p{11.9cm}|}

\hline
\rowcolor{gray!33} Value & Meaning \\ \hline
\endhead
\texttt{ok} & No executor is suspected. \\ \hline
\texttt{suspect} & At least one executor is suspected. \\ \hline
\texttt{unchecked} & None is suspected, but at least one refuses the question. \\ \hline
\texttt{disabled} & The check is off (\texttt{check-seconds} is 0). \\ \hline
\end{longtable}
}

\msubsection{enum}{LossCause}
Why an executor is suspected.
 Encoded as a JSON string.
{\small
\begin{longtable}{|p{3.9cm}|p{11.9cm}|}

\hline
\rowcolor{gray!33} Value & Meaning \\ \hline
\endhead
\texttt{notRunning} & It answers that it does not run the work item. \\ \hline
\texttt{unreachable} & It does not answer. \\ \hline
\texttt{unregistered} & It is no longer in the ServiceRegistry. \\ \hline
\end{longtable}
}

\msubsection{enum}{LossOutcome}
What follows for the step of a lost work item.
 Encoded as a JSON string.
{\small
\begin{longtable}{|p{3.9cm}|p{11.9cm}|}

\hline
\rowcolor{gray!33} Value & Meaning \\ \hline
\endhead
\texttt{nextExecutor} & It was not bound: the next executor takes it, without an extension. \\ \hline
\texttt{rebind} & x-rebind: another executor runs the step. \\ \hline
\texttt{retry} & x-retry: the same executor runs it after the backoff. \\ \hline
\texttt{sessionFailed} & S-FailInvoke: the retry policy allows no further attempt. \\ \hline
\end{longtable}
}

\msubsection{struct}{SuspectedExecutor}
An executor that did not confirm work items it holds.

{\small
\begin{longtable}{|p{3.7cm}|p{4.2cm}|p{1.6cm}|p{6.3cm}|}

\hline
\rowcolor{gray!33} Field & Type & Mandatory & Description \\ \hline
\endhead
\texttt{executor} & \pref{SystemName} & yes & The executor. \\ \hline
\texttt{cause} & \mref{LossCause} & yes & Why it is suspected. \\ \hline
\texttt{since} & \pref{DateTime} & yes & When its first work item was suspected. \\ \hline
\texttt{workItemCount} & \pref{Integer} & yes & Suspected work items of this executor. \\ \hline
\texttt{workItems} & \pref{List}$<$\mref{SuspectedWorkItem}$>$ & yes & The work items; at most 100 are listed over all executors. \\ \hline
\end{longtable}
}

\msubsection{struct}{SuspectedWorkItem}
A suspected work item.

{\small
\begin{longtable}{|p{3.7cm}|p{4.2cm}|p{1.6cm}|p{6.3cm}|}

\hline
\rowcolor{gray!33} Field & Type & Mandatory & Description \\ \hline
\endhead
\texttt{workItemId} & \pref{WorkItemId} & yes & The work item. \\ \hline
\texttt{sessionId} & \pref{SessionId} & yes & Its session. \\ \hline
\texttt{action} & \pref{Name} & yes & Its action. \\ \hline
\texttt{step} & \pref{Name} & yes & Its step. \\ \hline
\texttt{attempt} & \pref{Integer} & yes & Attempt number of the step. \\ \hline
\texttt{bound} & \pref{Boolean} & yes & Whether the executor reported BOUND, after which it uses the equipment. \\ \hline
\texttt{lastSeen} & \pref{DateTime} & yes & When the executor last confirmed or reported the work item. \\ \hline
\end{longtable}
}

\msubsection{struct}{WorkItemLost}
A work item that the WorkflowManager gave up because its executor is gone.

{\small
\begin{longtable}{|p{3.7cm}|p{4.2cm}|p{1.6cm}|p{6.3cm}|}

\hline
\rowcolor{gray!33} Field & Type & Mandatory & Description \\ \hline
\endhead
\texttt{workItemId} & \pref{WorkItemId} & yes & The work item. \\ \hline
\texttt{sessionId} & \pref{SessionId} & yes & Its session. \\ \hline
\texttt{recipe} & \pref{String} & yes & Recipe as \texttt{name@version}. \\ \hline
\texttt{action} & \pref{Name} & yes & Its action. \\ \hline
\texttt{step} & \pref{Name} & yes & Its step. \\ \hline
\texttt{attempt} & \pref{Integer} & yes & Attempt number of the step. \\ \hline
\texttt{executor} & \pref{SystemName} & yes & The executor that held it. \\ \hline
\texttt{cause} & \mref{LossCause} & yes & Why it was given up. \\ \hline
\texttt{detail} & \pref{String} & yes & The cause in words. \\ \hline
\texttt{bound} & \pref{Boolean} & yes & Whether it was bound. \\ \hline
\texttt{lastSeen} & \pref{DateTime} & yes & When the executor last confirmed or reported it. \\ \hline
\texttt{suspected\allowbreak{}Since} & \pref{DateTime} & yes & When the first check failed. \\ \hline
\texttt{lostAt} & \pref{DateTime} & yes & When it was given up. \\ \hline
\texttt{detection\allowbreak{}Seconds} & \pref{Duration} & yes & \texttt{lostAt} minus \texttt{lastSeen}: the longest time the loss can have gone unnoticed. \\ \hline
\texttt{equipment\allowbreak{}May\allowbreak{}Be\allowbreak{}In\allowbreak{}Use} & \pref{Boolean} & yes & True for a bound work item: an executor that is only cut off may still be running the step. \\ \hline
\texttt{outcome} & \mref{LossOutcome} & yes & What follows for the step. \\ \hline
\texttt{workRepeated} & \pref{Boolean} & yes & Whether a bound step is run again, so that its work may be done twice. \\ \hline
\end{longtable}
}


\subsection{Primitives}
\label{sec:model:primitives}
The abstract primitive types of the SD are represented in JSON~\cite{rfc8259} as follows.
\subsubsection*{alias \pdef{Boolean}}
JSON \texttt{true} or \texttt{false}.

\subsubsection*{alias \pdef{DateTime}}
JSON string in the RFC 3339 format~\cite{rfc3339}, e.g. \texttt{2026-\allowbreak{}11-\allowbreak{}25T10:15:30Z}.

\subsubsection*{alias \pdef{Duration}}
JSON number of seconds.

\subsubsection*{alias \pdef{Integer}}
JSON number without fraction.

\subsubsection*{alias \pdef{List}}
JSON array.

\subsubsection*{alias \pdef{Name}}
JSON string matching \texttt{\^{}[A-Za-z][A-Za-z0-9\_-]\{0,62\}\$}.

\subsubsection*{alias \pdef{SessionId}}
JSON number (64-bit integer); a path segment when used in a path.

\subsubsection*{alias \pdef{String}}
JSON string (UTF-8).

\subsubsection*{alias \pdef{SystemName}}
JSON string.

\subsubsection*{alias \pdef{WorkItemId}}
JSON string holding a UUID; a path segment when used in a path.


\section{Security}
\label{sec:security}
The connection uses TLS~\cite{rfc8446} with mutual authentication; the broker takes the user from the certificate (SASL
EXTERNAL). Its access control list must allow WorkflowManager to publish
\texttt{\{base\}.\#} (Section~\ref{sec:operations}), which covers the two types. The offer is registered with
interface template \texttt{generic\_\allowbreak{}amqps} and policy \texttt{CERT\_AUTH}. Nothing is published, and no event is kept,
while \texttt{sysmonitor.\allowbreak{}enabled} is false.

\section{Implementations}
\label{sec:implementations}
\begin{itemize}
\item \textbf{Java (reference).} Class \texttt{Executor\allowbreak{}Observations} of \texttt{workflow-\allowbreak{}manager}; the common library publishes
it (\texttt{Sys\allowbreak{}Monitor\allowbreak{}Source}, \texttt{Sys\allowbreak{}Monitor\allowbreak{}Publisher}).
\item \textbf{Go and Rust.} They have no workflow systems and so do not publish these types: the workflow systems are
\emph{not} interoperable across the three implementations. By inspection of their sources, the SysMonitor systems in Go and
Rust accept any observation type, so they can store and subscribe to these; this has not been tested with them.
\end{itemize}

\newpage
\bibliographystyle{IEEEtran}
\bibliography{bibliography}

\newpage

\section{Revision History}
\subsection{Amendments}

\noindent\begin{tabularx}{\textwidth}{| p{1cm} | p{2.2cm} | p{1.6cm} | X | p{3.6cm} |} \hline
\rowcolor{gray!33} No. & Date & Version & Subject of Amendments & Author \\ \hline
1 & 2026-11-25 & \arrowversion & Initial edition: interface design description of the WorkflowManager's sysMonitor observation types over AMQP/TLS/JSON, written from the Java reference implementation (ah5-Workflow-java). & Documentation maintenance \\ \hline
\end{tabularx}

\subsection{Quality Assurance}

\noindent\begin{tabularx}{\textwidth}{| p{1cm} | p{2.2cm} | p{1.6cm} | X |} \hline
\rowcolor{gray!33} No. & Date & Version & Approved by \\ \hline
1 & & \arrowversion & Pending review \\ \hline
\end{tabularx}

\end{document}
